\documentclass[11pt]{article}

\usepackage[margin=1in]{geometry}
\usepackage{amsmath,amssymb}
\usepackage[T1]{fontenc}
\usepackage{lmodern}
\usepackage{microtype}
\usepackage{xcolor}
\usepackage{listings}
\usepackage[hidelinks]{hyperref}
\usepackage{parskip}

\lstset{
  basicstyle=\ttfamily\small,
  frame=single,
  columns=fullflexible,
  breaklines=true
}

\title{An Intuitive Introduction to Lambda Calculus}
\author{}
\date{}

\begin{document}

\maketitle

Lambda calculus is easiest to understand as a \textbf{tiny programming
language where everything is built out of functions}.

Despite having almost nothing in it, lambda calculus is powerful enough
to express essentially any computation a normal programming language
can perform. It is one of the theoretical foundations of functional
programming and computer science.

\section{Start with an Ordinary Function}

Suppose you have
\[
f(x)=x+1.
\]

Lambda calculus asks: why bother giving the function a name such as
\(f\)? We can simply write the function itself:
\[
\lambda x.\;x+1.
\]

Read this as:

\begin{quote}
``A function that takes \(x\) and returns \(x+1\).''
\end{quote}

The symbol \(\lambda\) essentially means ``here comes an anonymous
function.''

In Python, the corresponding idea is:

\begin{lstlisting}[language=Python]
lambda x: x + 1
\end{lstlisting}

So the notation is not as exotic as it first looks.

\section{Applying a Function}

If we want to apply
\[
\lambda x.\;x+1
\]
to \(5\), we write
\[
(\lambda x.\;x+1)\;5.
\]

The intuitive rule is extremely simple:

\begin{quote}
Replace \(x\) with \(5\).
\end{quote}

Thus
\[
(\lambda x.\;x+1)\;5
\]
becomes
\[
5+1,
\]
and therefore
\[
6.
\]

This replacement process is called \textbf{beta reduction}.

\section{The Surprising Part: Pure Lambda Calculus Has Almost Nothing Else}

Strictly speaking, pure lambda calculus does not even start with
numbers, addition, \texttt{if} statements, loops, or Boolean values.

It has only three kinds of expressions:

\[
\boxed{x}
\]
a variable,

\[
\boxed{\lambda x.M}
\]
a function, and

\[
\boxed{M\;N}
\]
the application of one expression to another.

That is essentially the entire language.

The astonishing result is that \textbf{this is enough to represent
arbitrary computation}.

\section{Functions Can Consume Functions}

Consider
\[
\lambda f.\lambda x.f\;x.
\]

It means:

\begin{quote}
Take a function \(f\). Then take a value \(x\). Apply \(f\) to \(x\).
\end{quote}

For example, imagine
\[
f=\lambda y.y+1.
\]

Then
\[
(\lambda f.\lambda x.f\;x)\;f\;5
\]
eventually gives
\[
f(5)=6.
\]

This ability to treat functions just like ordinary values is the
central idea.

You can \textbf{pass a function into another function, return a
function from a function, and even use functions to represent data}.

\section{Numbers Can Be Functions}

Here is where lambda calculus becomes especially clever. It can
represent the number \(2\) as
\[
\lambda f.\lambda x.f(f(x)).
\]

Why would that mean \(2\)?

Think of a number as answering the question:

\begin{quote}
\textbf{How many times should I perform an operation?}
\end{quote}

Thus
\[
0=\lambda f.\lambda x.x
\]
means apply \(f\) zero times,

\[
1=\lambda f.\lambda x.f(x)
\]
means apply it once,

\[
2=\lambda f.\lambda x.f(f(x))
\]
means apply it twice, and

\[
3=\lambda f.\lambda x.f(f(f(x)))
\]
means apply it three times.

These are called \textbf{Church numerals}, after Alonzo Church, who
developed lambda calculus in the 1930s.

For instance, if
\[
f(x)=x+10,
\]
then Church numeral \(3\) gives
\[
f(f(f(0)))=30.
\]

The lambda expression itself does not know anything about the number
\(3\) in the conventional sense. ``Three'' is encoded as
\textbf{doing something three times}.

\section{Even TRUE and FALSE Can Be Functions}

Here is another elegant trick. Define
\[
\mathrm{TRUE}=\lambda x.\lambda y.x
\]
and
\[
\mathrm{FALSE}=\lambda x.\lambda y.y.
\]

Why?

\(\mathrm{TRUE}\) simply means:

\begin{quote}
Given two things, choose the first one.
\end{quote}

While \(\mathrm{FALSE}\) means:

\begin{quote}
Given two things, choose the second one.
\end{quote}

Now an \texttt{if} expression becomes almost trivial:
\[
\mathrm{IF}=\lambda b.\lambda x.\lambda y.b\;x\;y.
\]

Suppose
\[
\mathrm{IF}\;\mathrm{TRUE}\;A\;B.
\]

Since \(\mathrm{TRUE}\) chooses its first argument,
\[
\mathrm{TRUE}\;A\;B \longrightarrow A.
\]

Similarly,
\[
\mathrm{IF}\;\mathrm{FALSE}\;A\;B \longrightarrow B.
\]

So we have effectively constructed Boolean logic using
\textbf{nothing but functions}.

\section{The Big Idea}

Lambda calculus is not primarily interesting because it is a convenient
way to write programs. It is interesting because it asks a much deeper
question:

\begin{quote}
\textbf{What are the minimum ingredients necessary for computation?}
\end{quote}

It gives a remarkable answer:
\[
\boxed{\text{functions}+\text{variables}+\text{function application}}.
\]

These are enough.

Numbers can be represented by functions. Booleans can be represented by
functions. Data structures can be represented by functions. Recursion
can even be constructed from functions.

That is why lambda calculus can initially feel like someone saying
``everything is a function'' and then taking that idea absurdly
seriously---until you discover that it actually works.

\section{Where to Go Next}

After this intuitive introduction, the next important concepts are
\textbf{beta reduction and variable substitution}, especially the
distinction between \textbf{free and bound variables}. Those ideas form
the basis for working through lambda-calculus expressions formally.

\end{document}
